% Shared LaTeX preamble for homework/*/main.tex

% --- Base layout ---
\usepackage[a4paper, top=2.5cm, bottom=2.5cm, left=2cm, right=2cm]{geometry}
\usepackage{fontspec}
\usepackage[UTF8]{ctex}%latex支持中文宏包
\usepackage{amsmath, amssymb, amsthm}
\usepackage{mathtools}
\usepackage[svgnames]{xcolor}
\usepackage{pgfplots}
\pgfplotsset{compat=1.18}
\usepackage[most]{tcolorbox}
\usepackage{enumitem}
% Modified: Change label from hyphen (-) to bullet point ($\bullet$) for better visibility
\setlist[itemize]{label=$\bullet$}
\setlist[enumerate]{label=(\arabic*)}


% --- Colors ---
\definecolor{primaryColor}{RGB}{46, 52, 64}
\definecolor{accentColor}{RGB}{94, 129, 172}
\definecolor{boxFill}{RGB}{236, 240, 241}
\definecolor{solLine}{RGB}{163, 190, 140}
\definecolor{expandColor}{RGB}{191, 97, 106}

% --- TikZ styles (DFA / NFA / PDA) ---
\usepackage{tikz}
\usetikzlibrary{automata, positioning, arrows.meta, bending, backgrounds}
\tikzset{
    dfa_state/.style={
        state,
        thick,
        draw=accentColor,
        fill=boxFill,
        text=primaryColor,
        minimum size=1.1cm
    },
    dfa_edge/.style={
        ->,
        >=stealth,
        thick,
        draw=primaryColor,
        auto
    },
    pda_node/.style={
        state,
        thick,
        draw=accentColor,
        fill=boxFill,
        text=primaryColor,
        minimum size=1.2cm
    },
    pda_edge/.style={
        ->,
        >=stealth,
        thick,
        draw=primaryColor,
        auto
    },
    rule_expand/.style={
        pda_edge,
        draw=expandColor,
        text=expandColor
    },
    rule_match/.style={
        pda_edge,
        draw=solLine,
        text=solLine!80!black
    }
}

% --- Header / footer ---
\usepackage{fancyhdr}
\pagestyle{fancy}
\setlength{\headheight}{13.6pt}%避免页眉高度警告
\fancyhf{}
\fancyhead[L]{\kaishu 中国科学院大学}
\fancyhead[C]{林俊杰2023K8009916012}
\fancyhead[R]{\kaishu 理论计算机}
\usepackage{lastpage}
\cfoot{\small 第 \thepage \ 页，共 \pageref{LastPage} 页}

% --- Problem box ---
\newtcolorbox[use counter=problemSeq]{problem}[1][]{
    enhanced,
    breakable,
    colback=boxFill,
    colframe=accentColor,
    coltitle=white,
    fonttitle=\bfseries\large,
    title={题目 ~\theproblemSeq \quad #1},
    boxrule=0.5mm,
    arc=3mm,
    drop shadow=black!15!white,
    attach boxed title to top left={xshift=0.5cm, yshift*=-3mm},
    boxed title style={colback=accentColor},
    % 关键修改：每当创建一个新问题盒子时，自动将步骤计数器重置为0
    code={\setcounter{stepcount}{0}}
}

% --- Solution 环境 (红色盒子样式，支持可选序号) ---
% 修改说明：
% 1. [1][] 表示接受一个可选参数，默认为空
% 2. title={...} 中使用 \ifstrempty 判断是否有参数，如果有则添加空格和参数
\newtcolorbox{solution}[1][]{
    enhanced, breakable,
    colback=red!3!white,        % 背景色：淡粉
    colframe=red!75!black,      % 边框色：深红
    coltitle=white,
    fonttitle=\bfseries\Large\itshape,
    title={Solution\ifstrempty{#1}{}{~#1}}, % 动态标题：Solution (1)
    boxrule=0.5mm,
    arc=3mm,
    drop shadow=black!15!white,
    attach boxed title to top left={xshift=0.5cm, yshift*=-3mm},
    boxed title style={colback=red!75!black},
    before skip=0.5em,          % 盒子前的间距
    after skip=2em,             % 盒子后的间距
    before upper={\setcounter{stepcount}{0}} % 每个解答开始时重置步骤计数
}

% --- Title ---
\newcommand{\makeCustomTitle}[1]{
    \begin{center}
        \vspace*{1cm}
        {\Huge \bfseries \color{primaryColor} 作业 #1}\\
        \vspace{0.5cm}
        {\Large \color{accentColor} \today}
        \vspace{1.2cm}
    \end{center}
}

% --- 自定义环境定义 ---

% --- 自定义计数器与环境 ---
\newcounter{stepcount}
\newcounter{casecount}[stepcount]
\newcounter{problemSeq} % 新增：专门用于 Problem 的独立计数器

% 【关键修改】挂载钩子：每当遇到 \section 命令（含 \section*），将 problemSeq 计数器重置为 0
\pretocmd{\section}{\setcounter{problemSeq}{0}}{}{}

% 2. 定义带自动编号的步骤环境
\newenvironment{proofstep}[1]
  {%
    \refstepcounter{stepcount}% 增加计数
    \par\vspace{1.5em}%
    \noindent\textbf{\large \thestepcount. #1}% 输出 "1. 标题"
    \par\vspace{0.5em}%
  }
  {\par}

% 3. 定义无编号的结论环境 (用于最后的结论，格式同步骤)
\newenvironment{proofresult}[1]
  {%
    \par\vspace{1.5em}%
    \noindent\textbf{\large #1}% 仅输出标题
    \par\vspace{0.5em}%
  }
  {\par}

% 4. 定义带自动编号的情形环境
\newenvironment{proofcase}[1]
  {%
    \refstepcounter{casecount}% 增加计数
    \par\vspace{1em}%
    \noindent\textbf{情形 \thecasecount：#1}% 输出 "情形 1：标题"
    \par\vspace{0.2em}%
  }
  {\par}